\section{Varying the threefold volume character}\label{sec:variation}

Moving away from the volume-preserving torus of the threefold gives a
test of circle exchange that the numerical index of a single curve cannot
detect. Let \(C\) have genus two, \(B=\operatorname{Pic}^1(C)\), and
\(\pi:C\times B\to B\). Choose a Poincar\'e line \(\mathcal L\) trivialized
over one point of \(C\), and put
\(\mathcal M=K_C\mathcal L^{-1}\), \(A=R\pi_*\mathcal L\).
Give the four lines \(\mathcal L,\mathcal M,\mathcal O,\mathcal O\)
weights \(t,\kappa/t,q,\kappa^{-1}/q\), in that order. Their product
remains one on the fivefold. Write \(\delta=1-\kappa\) and retain the
orientation of the ordered external lines in the modified pair index.

Let \(\mathcal I_2^{\rm orig}\) and \(\mathcal I_2^{\rm rot}\) denote
the degree-two family characters with these weights in the original
and rotated presentations. We prove Theorem~\ref{thm:picard-variation}
by computing both sides of the identity
\begin{equation}
 [\delta q^{-3} t^k]\,\mathcal I_2^{\rm orig}
 =[\delta q^{-3} t^k]\,\mathcal I_2^{\rm rot}
 =-2\bigl(1+2\,\mathbf 1_{2\mid k}\bigr)\psi^k(A)
 \quad\text{in }K^0(B),
 \label{eq:circle-volume-edge}
\end{equation}
for every \(k\geq1\), where \(\psi^k\) is the Adams operation. In the original presentation
one takes the \(q^{-3}\) term, varies \(\kappa\), and expands the
result at \(t=0\). After exchange, \(t\) counts pairs. One fixes \(k\),
takes the variation, and expands at \(q=0\). In the latter expansion
there are no powers below \(q^{-3}\) at this order. We give the two
calculations because the fixed loci that produce the same class
have different dimensions.

\subsection{The two minimal ribbons}

At the lowest Euler characteristic, \(n=-3\), a torus-fixed degree-two
pair is the structure sheaf of a ribbon in either normal direction.
Indeed its generic cyclic quotient of length two has monomial basis
\(1,x\) or \(1,y\). The corresponding ribbon in a line \(F\) has
pushforward \(\mathcal O_C\oplus F^{-1}\), of Euler characteristic
\(-3\). A zero divisor of the section or of the multiplication between
the two weight lines raises the Euler characteristic. Thus each
choice gives one fixed copy of \(B\).

These two ribbons also occur as limits of smooth covers. For a
genus-two splitting as in \cite[Proposition~8.1]{wuebben-geometric},
the degree-two cycles form two one-parameter families, one in each
normal line. Their nonzero members are smooth genus-four curves.
In a fiber coordinate \(z\) on either line, the family has equation
\(z^2=a s\), where \(s\) is a section of the square of that line
with two simple zeros. At \(a=0\) the central ribbon has equation
\(z^2=0\). The index calculation below keeps its scheme structure
and its equivariant deformation weights.

Write \(F=\mathcal L,\ p=t\) on the first ribbon and
\(F=\mathcal M,\ p=\kappa/t\) on the second. The equations for the
ribbon are the other normal coordinate and the square of the first.
Tensoring the corresponding normal lines with
\(\mathcal O_C\oplus(pF)^{-1}\) and using relative Serre duality gives
the virtual tangent \(H-\kappa H^\vee\), with polarization
\[
 H=p R\pi_*F+p^2 R\pi_*F^2,\qquad \operatorname{rk}H=1 .
\]
For this polarization the symmetrized inverse Euler class, the
external determinant and the sign of \((-q)^{-3}\) give
\[
 \kappa^{-1}
 \frac{\lambda_{-1}(\kappa^{-1}H)}{\lambda_{-1}(H)} ,
 \qquad
 \lambda_{-1}(V)=\sum_j(-1)^j\mathop{\bigwedge}\nolimits^j V .
\]
The ratio at \(\kappa=1\) is identically one, so varying \(p=\kappa/t\)
in both its arguments contributes zero. A root \(w\) of \(H\)
contributes \(-w/(1-w)\) to its first variation. The expansions for
\(w=tx\) and \(w=t^{-1}x\) have opposite signs in positive degree
and involve \(x^k\) and \(x^{-k}\), respectively. Using
\(R\pi_*\mathcal M=-A^\vee\), the resulting positive coefficients are
\begin{equation}
 -2\psi^k(A)+\mathbf1_{2\mid k}\,
       \psi^{k/2}(D_M^\vee-D_L),\qquad
 D_L=R\pi_*\mathcal L^2,\quad D_M=R\pi_*\mathcal M^2 .
 \label{eq:circle-unsimplified-edge}
\end{equation}
Here and below the second term is included only for even \(k\).

\subsection{The fixed pairs after circle exchange}

The rotated threefold is \(C\times\mathbb A^2\). Write
\(u=q,\ v=\kappa^{-1}/q,\ c=uv\), so \(1-c^{-1}=\delta\).
A fixed row is specified by effective divisors \(P,Q\) on \(C\).
The section vanishes on \(P\) in the inner weight line, and the
map to the outer line vanishes on \(Q\). The pushforward to \(C\) and
the counting degree are
\[
 E=\mathcal O(P)\oplus u^{-1}\mathcal O(P+Q),
 \qquad k=2\deg P+\deg Q .
\]
The transposed row exchanges \(u\) and \(v\). Thus the fixed loci are
products of symmetric powers of \(C\). The virtual tangent of a row
is the pushforward along \(C\) of the following character, obtained
by resolving the pair in the two plane directions:
\[
 E+c E^\vee-(1-u)(1-v)\operatorname{End}E .
\]
For a universal divisor \(D\) of length \(d\), let \(p_D\) be its
projection to \(S_d=\operatorname{Sym}^d C\). Set
\(U_D=p_{D*}\mathcal O_D\) and
\(T_D=p_{D*}\mathcal O_D(D)=T_{S_d}\). The divisor sequences give
the weight-one and weight-\(c\) terms
\(T_P+T_Q-c(U_P+U_Q)\). The unit in \(U_D\) splits by the normalized
trace for \(d>0\). Each nonempty divisor therefore supplies a factor
\(1-c^{-1}\). All remaining moving weights are nontrivial at \(c=1\)
for generic \(q\). When both divisors are nonempty the first
variation vanishes, including along their incidence locus.

It remains to divide by the one factor \(1-c^{-1}\) on a pure
component and set \(c=1\). Write \(U=U_D,\ T=T_D,\ J=\det U\), and
write \(N(F)\) for the norm to \(S_d\) of a line \(F\) restricted
to \(D\). Finite duality gives
\(\det T=J^{-1}N(K_C)^{-1}\). The residual fixed obstruction
contributes \(\lambda_{-1}((U/\mathbf1)^\vee)\).
The moving tangent and modified root on the two components are
\begin{equation}
\begin{array}{c|rrrr|l}
 &q^{-2}&q^{-1}&q&q^2&\text{modified root}\\ \hline
 P=0,\ Q=D&T-1&-T-1&U-1&-U-1&q^d J N(\mathcal L)\\
 P=D,\ Q=0&-1&T-1&-U-1&-1&q^d J^2 N(\mathcal L^2).
\end{array}
\label{eq:circle-pure-moving}
\end{equation}
For example, subtract the minimal tangent
\(-(u+v+u^2+v/u)\) from the virtual tangent above. The divisor
sequences give \(T(1+v/u-v)+Uu(1-u-v)\) for the first row and
\(T(1+u^{-1})-cU(1+u)\) for the second.

For clarity, the orientation in the last column can be computed
before restricting to a fixed component. Write
\(\mathcal D(F)=\det R\pi_*F\), with the projection now taken over
the pair family, and let \(\langle\, ,\,\rangle\) be the Deligne
pairing on the curve. The modified root, up to constant
vector-space lines, is
\[
 c^{k+1}\mathcal D(E)^{-1}
             \langle\det E,\mathcal L K_C^{-1}\rangle .
\]
Indeed the endomorphism term multiplied by \((1-u)(1-v)\) has
trivial determinant. For rank-two \(E\) and a line \(F\) one has
\(\mathcal D(EF)/\mathcal D(E)
=\langle\det E,F\rangle
 (\mathcal D(F)/\mathcal D(\mathcal O))^2\).
This identity and Serre duality
\(\mathcal D(E^\vee)=\mathcal D(EK_C)\),
\(\mathcal D(\mathcal L)=\mathcal D(\mathcal M)\)
give the stated root on squaring, with external character exponent
\((\chi(E\mathcal M)-\chi(E^\vee))/2=k+1\).
Restriction to the universal divisors gives the norm lines in
\eqref{eq:circle-pure-moving}. The counting sign \((-1)^k\) is
still to be included.

A summand \(q^a V\) of the virtual tangent with \(a>0\) and virtual
rank \(r\) has inverse Euler class with leading term
\((-1)^r q^{ar}\det V\) at \(q=0\). A negative weight has leading
term one. The two rows of \eqref{eq:circle-pure-moving} start at
\(q^{-3}\), since the exponent from the root and positive weights is
\(d+(d-1)+2(-d-1)=-3\) in the first and
\(d+(-d-1)+2(-1)=-3\) in the second. At \(q=\infty\) the extreme
exponent is \(3\) for each row, which gives the same lower bound
for their transposes. To extract the coefficient put
\[
 \Phi_d(F)=R\pi_{S_d*}\bigl(N(F)\lambda_{-1}(U/\mathbf1)\bigr),
 \qquad \pi_{S_d}:S_d\times B\longrightarrow B .
\]
Combining the leading determinants with
\(J\lambda_{-1}((U/\mathbf1)^\vee)
=(-1)^{d-1}\lambda_{-1}(U/\mathbf1)\) and the counting sign gives
\(-\Phi_d(\mathcal L)\) from each outer divisor, at \(k=d\).
The common divisor and its transpose, at \(k=2d\), give respectively
\(-\Phi_d(\mathcal L^2)\) and
\(-\Phi_d(\mathcal L^2K_C^{-1})\). The factor \(N(K_C)^{-1}\)
in the transpose comes from \(\det T\) at the opposite end of
the \(q\)-expansion.

\subsection{Adams operations and the Picard polarization}

These symmetric-product indices admit an integral evaluation,
\begin{equation}
 \Phi_d(F)=\psi^d(R\pi_*F).
 \label{eq:circle-symmetric-adams}
\end{equation}
To see it, let \(h:C^d\to S_d\) be the ordered cover, let
\(V=\mathbb C^d\) be the permutation representation, and put
\(V_0=V/\mathbf1\). Marking one point of a divisor gives
\(U=(h_*\mathcal O\otimes V)^{\mathfrak S_d}\).
Exterior powers of these covariants identify with the covariants of
the exterior powers. Locally at a divisor of multiplicities \(d_i\)
this is Solomon's exterior algebra theorem for the product
\(\prod_i\mathfrak S_{d_i}\) acting on the local coordinates of the
ordered points \cite{solomon}. Splitting off the constant vector
gives the same statement for \(U/\mathbf1\) and \(V_0\).
The norm pulls back to \(F^{\boxtimes d}\). Relative K\"unneth now
expresses the pushforward in \eqref{eq:circle-symmetric-adams}
as the alternating sum of the Schur operations corresponding to
\(\bigwedge^jV_0\), with the usual Koszul signs on the tensor powers
of \(R\pi_*F\). Their partitions are the hooks \((d-j,1^j)\).
The integral symmetric-function identity
\(p_d=\sum_{j=0}^{d-1}(-1)^j s_{(d-j,1^j)}\)
identifies the sum with the Adams operation. Together with
\(D_M^\vee=-R\pi_*(\mathcal L^2K_C^{-1})\), this gives
\eqref{eq:circle-unsimplified-edge} directly in the rotated
presentation.

There is one more simplification in even counting degree. On
\(S_2\), \(J\) is the trace-free line of the ordered double cover,
so \(J^2=\mathcal O(-i(C))\) for the doubled-point divisor
\(i:C\to S_2\). Its restriction \(i^*J=K_C\) is the conormal
of the diagonal on the ordered cover. Multiplication by a line
minus one lowers the dimension of the support filtration on the
smooth surface \(S_2\). Hence \((J-1)^3=0\) in \(K^0(S_2)\), and
the divisor sequence gives
\[
 4(J-1)=(J-1)(2+J+J^{-1})
        =(J^2-1)(1+J^{-1})=-i_*(1+K_C^{-1}).
\]
Tensor with \(N(\mathcal L)\) and push to \(B\). Since \(U/\mathbf1=J\)
and \(i^*N(\mathcal L)=\mathcal L^2\), equation
\eqref{eq:circle-symmetric-adams} gives
\(4\psi^2(A)=R\pi_*(\mathcal L^2(1+K_C^{-1}))\).
Applying \(\psi^{k/2}\) proves \eqref{eq:circle-volume-edge}
integrally.

Finally let \(\theta\) be the principal polarization, with sign
fixed by \(\pi_*(\gamma^2)/2=-\theta\) for the mixed K\"unneth
component \(\gamma\) of \(c_1(\mathcal L)\). The degree-one line has
Euler characteristic zero on \(C\), and the chosen normalization
removes the base component of its first Chern class. Riemann--Roch
therefore gives \(\operatorname{ch}(A)=-\theta\) in \(H^{\mathrm{even}}(B,\mathbb Q)\). The cohomological Chern
character of \eqref{eq:circle-volume-edge} is
\(2k(1+2\,\mathbf1_{2\mid k})\theta\). At \(k=2\) the two outer
components give \(4\theta\) and the common-divisor components give
\(8\theta\). Their sum \(12\theta\) measures the response of the
family to varying the threefold volume character. Restriction to
any one target retains only its rank, which is zero.
